\documentclass[11pt]{article}
\usepackage[margin=0.75in]{geometry}



\begin{document}


\title{Aim 2 manuscript outline}
\author{Justin G. Chitpin}
\date{March 04, 2024.\\Last updated: \today}

\maketitle

\section{Abstract and preamble}

Title:\\

\noindent Atomic elementary flux modes explain the steady state flow of metabolites in flux networks.\\

\newline

\noindent Abstract:\\

TBD.\\

\noindent Keywords:
\begin{itemize}
    \item Flux network
    \item Elementary flux mode
    \item Flow decomposition
    \item Markov chain
    \item Atom tracing
\end{itemize}

\section{Introduction}

\begin{enumerate}
    \item Intro paragraph on metabolic flux networks.
    \begin{itemize}
        \item Possible hook: We can measure many properties of a cell but metabolic flux estimates are a `final' frontier for high-throughput assays.
        \item Advancements in experimental (SITA plus MS/NMR) and computational methods (FBA, metabolic flux analysis [13C-MFA], statistical models of flux).
        \item These fluxes can be represented as a knowledge-driven network. Insert definition.
        \item Many existing highly curated genome-scale models (acronym is either GeMs, GSM, or GSMMs depending whom you ask). Hundreds if not possibly thousands of these models.
        \item Different network structures across organisms but highly connected.
        \item Many ways to analyze these networks (cite examples).
        \item We focus on a fundamental question of tracing the steady state flow of metabolites in flux networks.
        \item Tracing metabolite flows is not only biophysically interesting but useful for quantifying changes in metabolism across conditions.
    \end{itemize}
    \item EFM background paragraph.
    \begin{itemize}
        \item EFMs are a pathway-based approach to explain steady state metabolite flows.
        \item Define EFMs and emphasize mechanistic/biophysically-motivated definition.
        \item Two main problems in EFM literature are the combinatorial explosion of pathways in large-scale networks and the flux decomposition problem.
        \item Computational problem of enumerating millions of EFMs plus cumbersome downstream analysis.
        \item Ambiguous EFM decompositions leads to non-predictive EFMs.
        \item We recently addressed the flux decomposition problem for the special case of unimolecular reaction networks.
        \item We modelled steady state particle flux as a Markov chain and developed an algorithm to uniquely identify EFM weights. Applied this to a unimolecular sphingolipid network with X metabolites and Y reactions.
    \end{itemize}
    \item Metabolic graphs paragraph.
    \begin{itemize}
        \item Graph-theoretic methods are also used to understand metabolite flows.
        \item Motivated by the many existing/fast algorithms to analyze large-scale networks.
        \item Used to identify metabolic hubs and characterize global/local network properties.
        \item Many existing methods to convert a metabolic network into a graph, although there are challenges preserving reaction directionality, steady state fluxes, and reaction stoichiometry. (Ex. Compound graphs, hypergraphs, reaction graphs, bipartite graphs.)
    \end{itemize}
    \item Flux-dependent graph paper paragraph.
    \begin{itemize}
        \item One recent paper to solve the problem of tracing metabolite flows by converting networks into a special type of reaction graph.
        \item Reaction graph where nodes are reactions and directed edges are the normalized probability of reaction $i$ producing metabolites $K$ that is subsequently consumed in reaction $j$.
        \item Many nice properties of this type of graph. Accounts for reaction directionality, stoichiometry, steady state flux conservation, and results in a graph that can be analyzed efficiently using existing algorithms.
        \item Major limitation is that pathways connecting reactions do not have clear biological meaning. In general, directed edges do not correspond to a single metabolite transferred between reactions but the aggregate of all metabolites shared between reaction pairs.
        \item The graph may satisfy flow conservation but a simple cycle or path does not explain the flow of a given metabolite through the network.
    \end{itemize}
    \item Directive paragraph.
    \begin{itemize}
        \item Longstanding problem of explaining the steady state flow of metabolites through flux networks. Two subproblems: (i) flux decomposition ambiguity in multispecies networks and (ii) pathways through reaction graphs cannot be interpreted as EFMs.
        \item We address these problems through a new metabolic network representation we call the atomic transition graph.
        \item The atomic transition graph is a directed graph that preserves reaction directionality, reaction stoichiometries, and steady state flux.
        \item It works by employing a state-of-the-art atom mapping algorithm to construct a graph that models the flow of an individual metabolite through the network.
        \item Notably, atomic EFMs are always linear and can be computed under an atomic version of the CHMC.
        \item We first introduce our framework for constructing and analyzing an atomic CHMC.
        \item Empirically demonstrate there are orders of magnitude fewer atomic EFMs than standard EFMs (which we should probably redefine as `molecular' EFMs) on 5 genome-scale networks, and orders of magnitude faster.
        \item Investigate structural properties of atomic EFMs: distributions, loops, classifications, etc.
        \item Differential atom metabolism and dominant flux hypothesis on that one set of HepG2 steady state fluxes?
    \end{itemize}
\end{enumerate}


\section{Figures/tables in (rough) order of appearance}

\begin{enumerate}
    \item Pipeline figure. (Completed)
    \item Table with summary statistics on the 5 datasets. (Completed)
    \item Benchmarking figures. (Completed)
    \item Structural analyses figure. (Histogram curves?)
    \item Jitter plots of differential metabolism. (Completed)
    \item HepG2 figure. (Dominant mass flow curves, networks?, pathway analysis? Bowtie visuals? Sankey diagrams?)

\end{enumerate}

\section{Results (MSB style)}

\subsection{Tracing atomic flows in genome-scale models/Pipeline overview}

\begin{itemize}
    \item Introduce algorithm/pipeline.
    \item Cite RXNMapper.
    \item 5 genome-scale models.
    \item Describe datasets. Ex:
    \begin{itemize}
        \item BIGG datasets and HepG2 dataset with a single set of steady state fluxes.
        \item HepG2 dataset constructed from the Human Metabolic Reaction (HMR 2.0) database.
        \item Steady-state fluxes estimated from time-resolved protein and metabolite abundance data and FBA.
        \item How were datasets picked? Smallest BiGG model and increasingly larger others.
    \end{itemize}
    \item Extracting SMILES strings (more details in Supplementary).
    \item Small comment that BiGG models and other GeMs do not include canonical or isomeric SMILES structures and these must be manually scraped. Also some disagreement on chemical formula for some metabolite species (I picked the SMILES structure linked to the BiGG database that matched where possible or else dropped that metabolite.)
    \item Mention convenience functions to check/correct the metabolic model inputs (unimolecular source/sink fluxes, integer stoichiometries, no pseudometabolites, etc.). Describe details in supplementary.
\end{itemize}

\subsection{Atomic EFMs are computationally tractable on genome-scale metabolic models/Benchmarking atomic-versus-standard EFMs}

\begin{itemize}
    \item Describe computational benefits of atomic-versus-molecular EFMs.
    \item Table of 5 model summary statistics, number of carbon/nitrogen-versus-molecular EFMs and benchmarking figures.
    \item Total number of atomic EFMs and standard EFMs.
    \item Total run time of carbon/nitrogen EFMs versus standard EFMs.
    \item 2-4 orders of magnitude fewer atomic EFMs and 2-4 orders of magnitude faster computational time.
    \item Log-linear increase in single-threaded time as a function of number of atomic EFMs.
    \item Log-linear makes sense given the combinatorial increase in atomic EFMs as the CHMC tree grows. Corroborated by math/CS papers studying EFMs  (ex. minimal cut sets paper by Vicente Acuña).
    \item Insert possible supplementary figure showing near perfect log-linear correlation between normalized number of atomic CHMC states and number of atomic EFMs.
    \item Atomic EFMs can be parallelized across each atom within a source metabolite while the majority of molecular EFM enumeration methods cannot.
    \begin{itemize}
        \item EFMlrs is possibly the only/latest parallelized molecular EFM enumeration tool.
        \item JCVI-syn3A dataset: 304 metabolites, 388 reactions. 12,051,382,513 molecular EFMs.
        \item 5 days with 960 threads.
    \end{itemize}
    \item Consider re-running all atomic CHMCs or a single (big) one with increasing number of threads to assess speedup as a function of number of threads.
    \item Possibly include an example multispecies EFM to explain why there are fewer atomic EFMs and why interpreting atomic EFMs is easier by focusing on nutrient sources.
    \item Caveats on oxygen/hydrogen/phosphorus EFMs with RXNMapper limitations. Also how oxygen/hydrogens are so ubiquitous tracing those pathways may not be important. Phosphorus tends to be conserved in ATP and its derivatives.
\end{itemize}

\subsection{A taxonomy of atomic EFMs}

\begin{itemize}
    \item High level understanding of atomic movements through a network.
    \item Source-to-sink EFMs explain how nutrient sources are remodelled.
    \item Internally looped EFMs explain how nutrient sources remain trapped within the network. Certain metabolites (and especially lipids) may exhibit high remodelling flux.
    \item We further hypothesized that some EFMs may involve an atom revisiting a metabolite at a different position. An example would be pyruvate undergoing the TCA cycle twice. We call these `metabolite loops'.
    \item Insert classification barplots.
    \item Discuss classification differences across datasets. Majority of atomic EFMs are source-to-sink with/without revisitations.
    \item I should find a small atomic EFM from each category and visualize the chemical structures with the highlighted carbon of interest.
\end{itemize}

\subsection{Structural analysis of atomic EFMs}

\subsubsection{Properties of atomic EFMs}

\begin{itemize}
    \item The number and type of atomic EFMs depend on the network structure, so studying structural properties of EFMs can reveal interesting properties of the metabolic network.
    \item Distribution of carbon/nitrogen EFM lengths.
    \begin{itemize}
        \item Unimodal carbon EFM distribution for the BiGG datasets.
        \item Unimodal carbon EFM distribution for bigger HepG2 dataset.
        \item Much broader/flatter nitrogen EFM distributions across 5 datasets.
        \item On average, HepG2 carbon EFMs transit 99 metabolites/atom states or 100 reactions.
        \item Caution that 99 states does not necessarily imply 99 distinct metabolites.
    \end{itemize}
    \item Distribution of metabolite loops weighted by the number of occurrences.
    \begin{itemize}
        \item An atom revisiting the same metabolite twice results in a single count.
        \item An atom revisiting the same metabolite thrice results in a double count.
    \end{itemize}
    \item Qualitatively*, I'm observing similar distribution shapes between carbon EFM lengths and carbon metabolite loops.
    \item Possible barplot(s) showing the breakdown of metabolite loops by internal metabolite for a specific dataset.
    \begin{itemize}
        \item Metabolites revisited the most are TCA intermediates and specific amino acids. The NAD+ metabolite is possibly explained by salvage pathways involving NAM, NA, NR and NMN.
    \end{itemize}
\end{itemize}


\subsubsection{Differential atom metabolism}

\begin{itemize}
    \item It is well-understood that nutrients/metabolites are broken down during metabolism because the majority of metabolic reactions are multispecies reactions.
    \item To the best of my knowledge, the question of quantifying the extent of differential atom remodelling has not been addressed in literature.
    \item We computed the log ratio of source metabolite carbons with the greatest and smallest number of EFMs.
    \item Data shows that specific metabolite carbons may undergo over 4 orders of magnitude more EFMs than other source metabolite carbons.
    \item Point out that the statistics for branched chain amino acids are particularly large given both carbons in glycine are equivalently metabolized.
    \item TCA intermediates like malate also exhibit fairly large but consistent levels of differential metabolism across the 5 datasets.
    \begin{itemize}
        \item Is this coincidence or does this point to some type of conservation in reactions involving malate and scaling across the various sized networks?
        \item Red blood cells do not contain mitochondria and that (likely?) explains the zero-valued differential statistic.
    \end{itemize}
    \item The statistics are large for threonine, arginine, phenylalanine, and methionine in the HepG2 dataset. Perhaps this is because that particular model includes a lot more metabolites/reactions involving those particular amino acids compared to the other datasets?
\end{itemize}


\subsection{Something on glutamine remodelling in HepG2 liver cancer cell line}

\begin{itemize}
    \item We next turn to...
    \item Dominant pathway flux hypothesis:
    \begin{itemize}
        \item Explain idea of questioning the distribution of pathway fluxes.
        \item Question can be addressed by studying the distribution of atomic EFM weights. But for a single set of steady state vectors...
        \item Cumulative atomic mass by highest ranking atomic EFMs plots (glucose and glutamine; rest in supplementary).
    \end{itemize}
    \item Sankey diagrams of top weighted glutamine carbon EFM bins to show source glutamine metabolism?
    \begin{itemize}
        \item Top 10 atomic EFMs explain 63.3\% of the total glutamine carbon mass flow. (Clarification: these 10 atomic EFMs may come from the one of five carbons in glutamine. The denominator is the total carbon mass flow of glutamine)
        \begin{itemize}
            \item Note: `percentage_explained_source_flux_by_efms(CHMCAtomicSummary)`
        \end{itemize}
        \item Maybe a Sankey diagram of the next 11-20 or maybe 41-50 glutamine carbon EFMs?
        \item Possibly compare the number of metabolites/reactions in the two Sankey diagrams?
        \item Assign more traditional metabolic pathways/compartments to the Sankey diagrams? Top 10 glutamine carbon EFMs clearly correspond to the malate-asparate shuttle for example.
        \item Note that EFMs with progressively smaller weights tend to contain more metabolites/reactions because the probabilities of a longer pathway will generally be smaller than shorter ones.
    \end{itemize}
    \item Bowtie visualizations to show how EFMs can pass through metabolites of interest?
    \begin{itemize}
        \item Bowtie center is the (internal) metabolite/atom state of interest. Left and right bows correspond to the source and sink metabolite/atom states that pass through the center. Possibly use line width/colour to denote the EFM pathway flux through source/sink pairs.
        \item Internal metabolites (predominantly amino acids) taken from the HepG2 PNAS paper.
    \end{itemize}
    \item Networks?
    \item Comparison between atomic EFMs and KEGG pathway/subsystems?
\end{itemize}


\section{Discussion (MSB version of a conclusion)}

\begin{itemize}
    \item Key points:
    \begin{itemize}
        \item Introduced concept of atomic EFMs. They can be enumerated and their weights computed by the CHMC and atom mapping algorithms.
        \item Demonstrated empirically that atomic EFMs are computationally tractable in genome-scale metabolic models versus their molecular EFM counterparts.
        \item Characterized structural properties of carbon and nitrogen EFMs in 5 datasets and quantified differential atom metabolism.
        \item Final points regarding HepG2 network and analyses involving the single set of steady state fluxes.
    \end{itemize}
    \item Discussion talking points:
    \begin{itemize}
        \item Applying our algorithm to larger datasets with more computational resources. Or perhaps compressing linear pathways to reduce the state space?
        \item Similarities with carbon flow paths and atomic carbon flow paths (CFPs and aCFPs) and atomic EFMs. Our work generalizes that to any atom type and solves the flux decomposition problem. https://genomebiology.biomedcentral.com/articles/10.1186/gb-2011-12-5-r49
        \item Comment that one can identify all atomic EFMs that produce a given sink metabolite by multiplying the stoichiometry matrix by -1 to construct reversed atomic CHMCs.
    \end{itemize}
    \item Further musings to possibly mention?
    \begin{itemize}
        \item Relationship between our work and C13 metabolic flux analysis. Both use atom mapping and network stoichiometry to model the movement of atoms. MFA yields isotopic distributions as a function of external and internal fluxes. This raises an interesting question about \textit{in silico} EFM analysis with MFA. One could ask how distributions of EFMs vary across different sets of fluxes that produce similar isotopic distributions. Perhaps there is much uncertainty regarding the fluxes but the atomic EFMs are similar or with the same rank.
        \item Can we go from isotopic distributions directly to atomic EFM weights? Or can we go from atomic EFM weight distributions directly to metabolic flux distributions?
    \end{itemize}
\end{itemize}


\section{Methods}

\subsection{Pre-processing genome-scale models}

\begin{itemize}
    \item Dataset publications/citations.
    \item How many pseudo-metabolites removed?
    \item How many non-integer reactions removed/converted into unimolecular sources/sinks?
    \item How many reaction SMILES strings removed because of RXNMapper character limit?
    \item Explicitly representing reversible reactions as pairs of irreversible ones.
    \item Ensuring source/sink fluxes to stoichiometries of one.
    \item Splitting intracellular carbon dioxide into sources/sinks.
    \item All pre-processing steps above are handled by my software package.
\end{itemize}

\subsection{Extracting SMILES strings}

\begin{itemize}
    \item HepG2 metabolite SMILES taken from the PubChem database. Isomeric SMILES used where applicable.
    \item BiGG database metabolite SMILES taken from either PubChem or MetaNetX database which is referenced on the BiGG website.
    \item The same SMILES were used for the same metabolite found within different compartments.
    \item Metabolite SMILES were not recorded when there was ambiguity regarding the compound structure that could not be resolved through another chemical compound database.
\end{itemize}

\subsection{Atomic-versus-standard EFM benchmarking}

\begin{itemize}
    \item Hardware specs.
    \item MATLAB/Julia versions.
    \item Note we precompile Julia functions separately before running benchmarks.
    \item Benchmarks do not include startup time to open a Julia or MATLAB session.
\end{itemize}

\subsection{Atomic compound network visualization (if any we show any networks)}

\begin{itemize}
    \item Gephi version.
    \item Yifan Hu layout.
\end{itemize}


\section{Constructing atomic CHMCs}

\begin{itemize}
    \item Supplementary figure describing how atomic fluxes are assigned based on reaction fluxes.
    \item Atom mappings for reversible reactions were taken from atom mappings of the forward reaction for consistency.
\end{itemize}

\section{Solving for atomic EFM probabilities}

\begin{itemize}
    \item Redefine notation for solving for EFM probabilities.
    \item Computational issues solving for the stationary distribution of a very large CHMC.
    \item This involves solving for the eigenvector of the transition matrix.
    \item The eigenvector can be solved by direct or iterative methods.
    \item Pros/cons for direct/iterative solvers.
    \item As a rule of thumb, iterative solvers tend to perform faster on sparse matrices with regular structures.
    \item Iterative methods have tolerance parameters and we found that methods such as the Arnoldi method from `Arnoldi.jl' outputed negative values.
    \item Some issues obtaining NaNs or negative values using iterative methods such as GMRES from `IterativeSolvers.jl'.
    \item For direct or iterative methods, we found that repeatedly running the function would result in a correct eigenvector with no NaN/negative eigenvector coefficients.
    \item We've found that the package `LinearSolve.jl' worked well using default parameters and repeatedly running it on large matrices resulted in identical eigenvalues.
\end{itemize}

\section{Data availability}

\begin{itemize}
    \item BiGG database link.
    \item Code repository link.
    \item Package repository link.
\end{itemize}


\section{Acknowledgements/Funding}

This research was supported by the Natural Sciences and Engineering Research Council of Canada (NSERC), Discovery grant RGPIN-2019-0660. J.G.C. was supported by an NSERC CREATE Matrix Metabolomics Scholarship and an NSERC Alexander Graham Bell Canada Graduate Scholarship.

\section{Author contributions}

Justin G. Chitpin: Conceptualization; data curation; methodology; software; formal analysis; investigation; visualization; writing - original draft; writing - review and editing.

Theodore J. Perkins: Conceptualization; methodology; supervision; funding acquisition; investigation; project administration; writing - review and editing.

\section{Disclosure of competing interests}

The authors declare no competing interests.

\end{document}
